problem:  
Let \( M^n \) be a compact, simply connected Riemannian manifold with **positive sectional curvature**, and let the identity component of its isometry group \(G = \mathrm{Isom}_0(M)\) contain a **semisimple subgroup** \( K \) acting with **cohomogeneity four**.  For a principal orbit \( K/H \subset M \), the normal 4‑plane slice carries the orthogonal \(H\)-representation  
\[
\rho: H \to \mathrm{SO}(4) \cong (\mathrm{SU}(2)_+\times \mathrm{SU}(2)_-)/\{\pm 1\}.
\]  
Write \(\pi_\pm\circ\rho : H \to \mathrm{SU}(2)_\pm\) for the two component projections.  

**Known examples** of positively curved manifolds with this type of symmetry (e.g. certain homogeneous spaces and biquotients) all satisfy that **exactly one** of these projections, \(\pi_+\circ\rho\) or \(\pi_-\circ\rho\), is trivial—i.e. \(\rho(H)\) lies in a conjugate of either \(\mathrm{SU}(2)_+\) or \(\mathrm{SU}(2)_-\), corresponding to a reduction to \(\mathrm{Sp}(1)\) or \(\mathrm{U}(2)\).

**Question:**  
Does positive sectional curvature *force* this one‑sidedness condition?  
Equivalently, can there exist a compact, positively curved manifold \(M^n\) with a semisimple cohomogeneity‑four action for which both projections  
\[
\pi_+\circ\rho,\; \pi_-\circ\rho
\]
are nontrivial—so that the slice isotropy \(\rho(H)\) couples the two \(\mathrm{SU}(2)\) factors intrinsically, lying in no conjugate of \(\mathrm{U}(2)\) or \(\mathrm{Sp}(1)\)?  

If such “mixed’’ slice representations cannot occur, positive curvature would impose an unexpected rigidity forcing isotropy‑slice interactions to decouple into purely left‑ or right‑handed actions; if they can, explicit constructions would open a new family of semisymmetric positively curved spaces with novel isotropy geometry.

Why it is a "good" problem:  
This version makes the representation‑theoretic condition **precise and nontrivial**.  It probes whether positive curvature enforces an intrinsic **chirality restriction** on \(\mathrm{SO}(4)\)‑slices, excluding actions where isotropy mixes the two \(\mathrm{SU}(2)\) factors.  

The problem unites  
- **Riemannian geometry** (curvature positivity and orbit geometry),  
- **transformation group theory** (low‑cohomogeneity actions), and  
- **representation theory** (coupling between \(\mathrm{SU}(2)_+\) and \(\mathrm{SU}(2)_-\)).  

A proof of such a rigidity would uncover a new structural phenomenon: curvature dictating the algebraic decomposition of isotropy actions. A counterexample would yield an entirely new geometry of positively curved manifolds, challenging the prevailing quaternionic/complex bias of known models.  
It is succinct, logically coherent, and conceptually rich—meeting the criteria for a “good’’ research‑level mathematical problem.